\chapter{Chapter05 — Threshold Circuits}
%%% generated by the blueprint pipeline from FABL.Chapter05.ThresholdCircuits
\section{Overview}

Declarations of the Lean module \texttt{FABL.Chapter05.ThresholdCircuits}.

\section{Declarations}

\begin{definition}[Real sums of linear thresholds]
\lean{FABL.IsRealSumOfLinearThresholds}
\statementsource{Analysis-of-Boolean-Functions-by-Ryan-ODonnell}{
  pdf-ffbe9c7b3306-p044-b005,
  pdf-ffbe9c7b3306-p135-b006
}
\label{FABL.IsRealSumOfLinearThresholds}
\leanok
\uses{FABL.BooleanFunction, FABL.BooleanFunction.toReal, FABL.IsLinearThreshold}
Let $n$ and $s$ be nonnegative integers, and let $f$ be a Boolean function on
$\{-1,1\}^n$ with values in $\{-1,1\}$. The function $f$ is a real sum of at most
$s$ linear threshold functions if there are linear threshold functions $g_1,\ldots,g_s$
and a real constant $c$ such that, at every input $x$, the real encoding of $f(x)$
equals $c+\sum_{i=1}^s g_i(x)$, with each $g_i(x)$ interpreted by its real encoding.
\end{definition}

\begin{definition}[Real-sum size bound]
\lean{FABL.HasRealSumOfLinearThresholdsSizeAtMost}
\statementsource{Analysis-of-Boolean-Functions-by-Ryan-ODonnell}{pdf-ffbe9c7b3306-p135-b006}
\label{FABL.HasRealSumOfLinearThresholdsSizeAtMost}
\leanok
\uses{FABL.BooleanFunction, FABL.IsRealSumOfLinearThresholds}
Let $n$ and $s$ be natural numbers, and let $f$ be a function from the $n$-dimensional
Boolean cube to $\{ -1,1 \}$. We say that $f$ has real-sum size at most $s$ if there
is a natural number $t\le s$ such that $f$ is a real sum of $t$ linear threshold
functions.
\end{definition}

\begin{definition}[Threshold of linear thresholds]
\lean{FABL.IsThresholdOfThresholds}
\statementsource{freund-schapire-hedge}{pdf-01f49de027c4-p010-b006}
\label{FABL.IsThresholdOfThresholds}
\leanok
\uses{FABL.BooleanFunction, FABL.IsLinearThreshold, FABL.signValue, FABL.thresholdSign}
Let $n$ and $s$ be nonnegative integers, and let $f$ be a Boolean function on
$\{-1,1\}^n$ with values in $\{-1,1\}$. The function $f$ is a threshold of $s$
linear threshold functions if there are linear threshold functions $g_1,\ldots,g_s$,
real coefficients $a_0,a_1,\ldots,a_s$, and, at every input $x$, $f(x)$ equals the sign
threshold of $a_0+\sum_{i=1}^s a_i g_i(x)$, where each $g_i(x)$ is viewed as a sign.
\end{definition}

\begin{definition}[Threshold-of-thresholds size bound]
\lean{FABL.HasThresholdOfThresholdsSizeAtMost}
\label{FABL.HasThresholdOfThresholdsSizeAtMost}
\leanok
\uses{FABL.BooleanFunction, FABL.IsThresholdOfThresholds}
For a Boolean function $f$ on the sign cube of dimension $n$ and a natural number $s$,
the function has a threshold-of-thresholds representation of size at most $s$ if there
is a natural number $t\le s$ such that $f$ is a threshold of thresholds of size $t$.
\end{definition}

\begin{definition}[Sparse polynomial threshold representation]
\lean{FABL.IsThresholdOfParities}
\statementsource{Analysis-of-Boolean-Functions-by-Ryan-ODonnell}{
  pdf-ffbe9c7b3306-p115-b009,
  pdf-ffbe9c7b3306-p116-b001
}
\label{FABL.IsThresholdOfParities}
\leanok
\uses{FABL.BooleanFunction, FABL.IsPolynomialThresholdRepresentation, FABL.polynomialSparsity}
Let $n$ and $s$ be nonnegative integers, and let $f$ be a Boolean function on
$\{-1,1\}^n$ with values in $\{-1,1\}$. There is a real-valued polynomial on the
Boolean cube that represents $f$ as a polynomial threshold function and has sparsity at
most $s$.
\end{definition}

\begin{definition}[Positive-coordinate count on a set]
\lean{FABL.positiveCoordinateCountOn}
\label{FABL.positiveCoordinateCountOn}
\leanok
For a finite set $S\subseteq\{1,\ldots,n\}$ and a sign vector $x\in\{-1,1\}^n$, the
positive-coordinate count on $S$ is the number of indices $i\in S$ for which $x_i=1$.
\end{definition}

\begin{theorem}[Positive coordinate count bound on a subset]
\lean{FABL.positiveCoordinateCountOn_le}
\label{FABL.positiveCoordinateCountOn_le}
\leanok
\uses{FABL.positiveCoordinateCountOn}
\difficulty{1}
Let $S$ be a subset of the coordinates $\in n$, and let $x\in\{-1,1\}^n$. The number
of coordinates in $S$ where $x$ equals $1$ is at most $|S|$.
\end{theorem}

\begin{theorem}[Sign sum on a coordinate subset]
\lean{FABL.sum_signValue_on_eq}
\label{FABL.sum_signValue_on_eq}
\leanok
\uses{FABL.positiveCoordinateCountOn, FABL.signValue, FABL.signValue_eq_ite}
\difficulty{3}
Let $n$ be a nonnegative integer, let $S\subseteq\{1,\ldots,n\}$, and let
$x\in\{-1,1\}^n$. The sum of the signs $x_i$ over $i\in S$ equals twice the number of
coordinates in $S$ where $x_i=1$, minus $|S|$.
\end{theorem}

\begin{definition}[Positive-coordinate count threshold gate]
\lean{FABL.countThresholdGate}
\label{FABL.countThresholdGate}
\leanok
\uses{FABL.BooleanFunction, FABL.positiveCoordinateCountOn}
For a finite coordinate set $S\subseteq\{1,\ldots,n\}$ and an integer threshold $k\ge
0$, define a sign-valued Boolean function on $\{−1,1\}^n$ that equals $1$ when at least
$k$ coordinates in $S$ are positive, and equals $−1$ otherwise.
\end{definition}

\begin{theorem}[Count threshold gates are linear thresholds]
\lean{FABL.countThresholdGate_isLinearThreshold}
\label{FABL.countThresholdGate_isLinearThreshold}
\leanok
\uses{FABL.IsLinearThreshold, FABL.countThresholdGate, FABL.positiveCoordinateCountOn, FABL.signValue, FABL.sum_signValue_on_eq, FABL.thresholdSign_of_neg, FABL.thresholdSign_of_nonneg}
\difficulty{5}
Let $n$ be a natural number, let $S\subseteq\{0,\ldots,n-1\}$, and let $k$ be a natural
number. On the sign cube, the function that outputs $1$ when at least $k$ coordinates
indexed by $S$ are positive and outputs $-1$ otherwise is a linear threshold function:
there are real coefficients and a real bias whose threshold sign gives this function at
every input.
\end{theorem}

\begin{definition}[Negative count threshold gate]
\lean{FABL.negCountThresholdGate}
\label{FABL.negCountThresholdGate}
\leanok
\uses{FABL.BooleanFunction, FABL.positiveCoordinateCountOn}
For a finite coordinate set $S\subseteq\{1,\ldots,n\}$ and an integer threshold $k\ge
0$, define a sign-valued Boolean function on $\{−1,1\}^n$ that equals $−1$ when at least
$k$ coordinates in $S$ are positive, and equals $1$ otherwise.
\end{definition}

\begin{theorem}[A negative-count threshold is linear]
\lean{FABL.negCountThresholdGate_isLinearThreshold}
\label{FABL.negCountThresholdGate_isLinearThreshold}
\leanok
\uses{FABL.IsLinearThreshold, FABL.negCountThresholdGate, FABL.positiveCoordinateCountOn, FABL.signValue, FABL.sum_signValue_on_eq, FABL.thresholdSign_of_neg, FABL.thresholdSign_of_nonneg}
\difficulty{5}
Let $n$ be a natural number, let $S$ be a subset of the coordinates, and let $k$ be a
natural number. Define a sign-valued gate to output $-1$ when at least $k$ coordinates
in $S$ are positive, and $1$ otherwise. This gate is a linear threshold function: there
are real weights and a real offset such that its value at each sign vector is the sign
of the resulting affine linear form.
\end{theorem}

\begin{definition}[Constant sign function]
\lean{FABL.constantThresholdGate}
\label{FABL.constantThresholdGate}
\leanok
\uses{FABL.BooleanFunction, FABL.Sign}
For any sign $b\in\{−1,1\}$, define the Boolean function on $\{−1,1\}^n$ that takes the
constant value $b$.
\end{definition}

\begin{theorem}[Constant functions are linear threshold functions]
\lean{FABL.constantThresholdGate_isLinearThreshold}
\label{FABL.constantThresholdGate_isLinearThreshold}
\leanok
\uses{FABL.IsLinearThreshold, FABL.Sign, FABL.constantThresholdGate}
\difficulty{4}
For any natural number $n$, every constant Boolean function on $\{-1,1\}^n$ is a linear
threshold function: there are a real threshold offset and real weights on the
coordinates such that the function's value at each input is the sign of the offset plus
the weighted sum of its input coordinates, with zero mapped to $1$.
\end{theorem}

\begin{definition}[Profile main gate]
\lean{FABL.profileMainGate}
\label{FABL.profileMainGate}
\leanok
\uses{FABL.BooleanFunction, FABL.Sign, FABL.constantThresholdGate, FABL.countThresholdGate, FABL.negCountThresholdGate}
Let $S\subseteq\{1,\ldots,n\}$ be finite, let $a\colon\mathbb{N}\to\{−1,1\}$ be a sign
sequence, and let $i\in\{0,\ldots,n-1\}$. Define a Boolean function on $\{−1,1\}^n$ as
follows: if $a(i+1)=a(i)$, it is constantly $1$; otherwise, if $a(i+1)=1$, it is the
gate that outputs $1$ when at least $i+1$ coordinates in $S$ are positive and $−1$
otherwise; if $a(i+1)=−1$, it outputs $−1$ when at least $i+1$ coordinates in $S$ are
positive and $1$ otherwise.
\end{definition}

\begin{definition}[Profile correction gate]
\lean{FABL.profileCorrectionGate}
\label{FABL.profileCorrectionGate}
\leanok
\uses{FABL.BooleanFunction, FABL.Sign, FABL.constantThresholdGate}
Let $a\colon\mathbb{N}\to\{−1,1\}$ be a sign sequence and let $i\in\{0,\ldots,n-1\}$.
Define a constant Boolean function on $\{−1,1\}^n$: its value is $−1$ if $a(i+1)=a(i)$,
and is $a(i+1)$ otherwise.
\end{definition}

\begin{theorem}[Linearity of the profile main gate]
\lean{FABL.profileMainGate_isLinearThreshold}
\label{FABL.profileMainGate_isLinearThreshold}
\leanok
\uses{FABL.IsLinearThreshold, FABL.Sign, FABL.constantThresholdGate_isLinearThreshold, FABL.countThresholdGate_isLinearThreshold, FABL.negCountThresholdGate_isLinearThreshold, FABL.profileMainGate}
\difficulty{4}
Let $n$ be a natural number, let $S$ be a subset of the $n$ coordinate indices, let
$a:\mathbb{N}\to\{−1,1\}$ be a sign sequence, and choose a coordinate index $i$. The
main gate that is constantly $1$ when $a(i+1)=a(i)$, and otherwise outputs $1$ or $-1$
according as the number of coordinates in $S$ taking value $1$ is at least $i+1$ or is
smaller, with the output direction chosen by whether $a(i+1)=1$, is a linear threshold
function.
\end{theorem}

\begin{theorem}[Linearity of a profile correction gate]
\lean{FABL.profileCorrectionGate_isLinearThreshold}
\label{FABL.profileCorrectionGate_isLinearThreshold}
\leanok
\uses{FABL.IsLinearThreshold, FABL.Sign, FABL.constantThresholdGate_isLinearThreshold, FABL.profileCorrectionGate}
\difficulty{3}
Let $n$ be a natural number, let $a:\mathbb{N}\to\{−1,1\}$ be a sign sequence, and
choose a coordinate index $i\in\{0,\ldots,n-1\}$. The gate that is constantly $-1$ when
$a(i+1)=a(i)$ and otherwise constantly $a(i+1)$ is a linear threshold function on the
$n$-dimensional sign cube.
\end{theorem}

\begin{theorem}[Value of a profile gate pair]
\lean{FABL.profileGate_pair_value}
\label{FABL.profileGate_pair_value}
\leanok
\uses{FABL.BooleanFunction.toReal, FABL.Sign, FABL.constantThresholdGate, FABL.countThresholdGate, FABL.negCountThresholdGate, FABL.positiveCoordinateCountOn, FABL.profileCorrectionGate, FABL.profileMainGate, FABL.signValue}
\difficulty{5}
Let $n$ be a natural number, let $S$ be a subset of the $n$ coordinate indices, let
$a:\mathbb{N}\to\{−1,1\}$ be a sign sequence, choose a coordinate index $i$, and let $x$
be a sign vector. The sum of the real-valued outputs of the corresponding main and
correction gates is $a(i+1)-a(i)$ if $i$ is less than the number of coordinates in $S$
where $x$ equals $1$, and is $0$ otherwise.
\end{theorem}

\begin{theorem}[Sum of profile gate pairs]
\lean{FABL.sum_profileGate_pair}
\label{FABL.sum_profileGate_pair}
\leanok
\uses{FABL.BooleanFunction.toReal, FABL.Sign, FABL.positiveCoordinateCountOn, FABL.positiveCoordinateCountOn_le, FABL.profileCorrectionGate, FABL.profileGate_pair_value, FABL.profileMainGate, FABL.signValue}
\difficulty{4}
Let $n$ be a nonnegative integer, let $S$ be a subset of the coordinate set
$\{1,\ldots,n\}$, let $a$ be a sequence of signs indexed by the nonnegative integers,
and let $x$ be a sign vector of length $n$. For each coordinate $i$, add the real-valued
outputs at $x$ of the associated profile main gate and profile correction gate. The sum
of these pairs over all coordinates is the sign value $a(c)-a(0)$, where $c$ is the
number of coordinates in $S$ at which $x$ equals $1$.
\end{theorem}

\begin{definition}[Profile decomposition gates]
\lean{FABL.profileDecompositionGates}
\label{FABL.profileDecompositionGates}
\leanok
\uses{FABL.BooleanFunction, FABL.Sign, FABL.profileCorrectionGate, FABL.profileMainGate}
Let $S\subseteq\{1,\ldots,n\}$ be finite and let $a\colon\mathbb{N}\to\{−1,1\}$ be a
sign sequence. Define a family of $2n$ Boolean functions indexed by pairs $(i,r)$ with
$0\le i<n$ and $r\in\{0,1\}$: for $r=0$ use the profile main gate at $i$, and for $r=1$
use the profile correction gate at $i$.
\end{definition}

\begin{theorem}[Sum of profile decomposition gates]
\lean{FABL.sum_profileDecompositionGates}
\label{FABL.sum_profileDecompositionGates}
\leanok
\uses{FABL.BooleanFunction.toReal, FABL.Sign, FABL.positiveCoordinateCountOn, FABL.profileDecompositionGates, FABL.signValue, FABL.sum_profileGate_pair}
\difficulty{3}
Let $n$ be a nonnegative integer, let $S$ be a subset of the coordinate set
$\{1,\ldots,n\}$, let $a$ be a sequence of signs indexed by the nonnegative integers,
and let $x$ be a sign vector of length $n$. The sum of the real-valued outputs of the
$2n$ profile decomposition gates associated with $S$ and $a$, evaluated at $x$, is the
sign value $a(c)-a(0)$, where $c$ is the number of coordinates in $S$ at which $x$
equals $1$.
\end{theorem}

\begin{theorem}[Linearity of the profile decomposition gates]
\lean{FABL.profileDecompositionGates_isLinearThreshold}
\label{FABL.profileDecompositionGates_isLinearThreshold}
\leanok
\uses{FABL.IsLinearThreshold, FABL.Sign, FABL.profileCorrectionGate_isLinearThreshold, FABL.profileDecompositionGates, FABL.profileMainGate_isLinearThreshold}
\difficulty{3}
Let $n$ be a natural number, let $S$ be a subset of the $n$ coordinate indices, and let
$a:\mathbb{N}\to\{−1,1\}$ be a sign sequence. For each coordinate $i$, form a main gate
that is constantly $1$ if $a(i+1)=a(i)$, and otherwise tests whether the number of
coordinates in $S$ taking value $1$ is at least $i+1$, with output $1$ when $a(i+1)=1$
and output $-1$ otherwise. Pair it with a correction gate that is constantly $-1$ if
$a(i+1)=a(i)$ and otherwise constantly $a(i+1)$. Every gate in this resulting family of
$2n$ gates is a linear threshold function.
\end{theorem}

\begin{theorem}[Count profiles as sums of linear thresholds]
\lean{FABL.isRealSumOfLinearThresholds_of_count_profile}
\label{FABL.isRealSumOfLinearThresholds_of_count_profile}
\leanok
\uses{FABL.BooleanFunction, FABL.IsRealSumOfLinearThresholds, FABL.Sign, FABL.positiveCoordinateCountOn, FABL.profileDecompositionGates, FABL.profileDecompositionGates_isLinearThreshold, FABL.signValue, FABL.sum_profileDecompositionGates}
\difficulty{3}
Let $n$ be a nonnegative integer, let $S\subseteq\{1,\ldots,n\}$, let
$a:\mathbb{N}\to\{-1,1\}$, and let $f:\{-1,1\}^n\to\{-1,1\}$. For $x\in\{-1,1\}^n$,
count the coordinates $i\in S$ with $x_i=1$. If $f(x)$ equals $a$ applied to this count
for every $x$, then the real-valued version of $f$ is a constant plus a sum of $2n$
linear threshold functions on the cube.
\end{theorem}

\begin{theorem}[Symmetric functions as sums of thresholds]
\lean{FABL.symmetric_hasRealSumOfLinearThresholdsSizeAtMost_two_mul}
\label{FABL.symmetric_hasRealSumOfLinearThresholdsSizeAtMost_two_mul}
\leanok
\uses{FABL.BooleanFunction, FABL.HasRealSumOfLinearThresholdsSizeAtMost, FABL.IsSymmetric, FABL.Sign, FABL.canonicalCountInput, FABL.isRealSumOfLinearThresholds_of_count_profile, FABL.isSymmetric_iff_eq_of_positiveCoordinateCount_eq, FABL.positiveCoordinateCount, FABL.positiveCoordinateCountOn, FABL.positiveCoordinateCount_canonicalCountInput, FABL.positiveCoordinateCount_le}
\difficulty{4}
Let $f$ be a function from the $n$-dimensional sign cube to $\{−1,1\}$ that is
invariant under every permutation of its input coordinates. Then there are an integer
$t\le 2n$, a real constant $c$, and $t$ linear threshold functions $g_1,\ldots,g_t$ on
the same cube such that, for every input $x$, the real-valued version of $f$ satisfies
$f(x)=c+\sum_{j=1}^t g_j(x)$, where each $g_j(x)$ is interpreted as its real sign
value.
\end{theorem}

\begin{theorem}[Threshold of linear thresholds]
\lean{FABL.isThresholdOfThresholds_of_isRealSumOfLinearThresholds}
\label{FABL.isThresholdOfThresholds_of_isRealSumOfLinearThresholds}
\leanok
\uses{FABL.BooleanFunction, FABL.BooleanFunction.toReal, FABL.IsRealSumOfLinearThresholds, FABL.IsThresholdOfThresholds}
\difficulty{4}
Let $n$ and $s$ be nonnegative integers, and let $f:\{\pm1\}^n\to\{\pm1\}$. Suppose
there are $s$ linear threshold functions $g_1,\ldots,g_s$ and a real constant $c$ such
that, for every input $x$, the real value of $f(x)$ is $c+\sum_{i=1}^s g_i(x)$, where
each sign is viewed as its real value. Then $f$ is a threshold of those $s$ linear
threshold functions: there are real weights and a real threshold offset such that $f(x)$
is the sign of the resulting weighted sum of $g_i(x)$, with zero assigned sign $+1$.
\end{theorem}

\begin{definition}[Negative-coordinate count]
\lean{FABL.negativeCoordinateCountOn}
\label{FABL.negativeCoordinateCountOn}
\leanok
For a finite coordinate set $S\subseteq\{1,\ldots,n\}$ and a point $x\in\{−1,1\}^n$,
define the negative-coordinate count on $S$ to be the number of coordinates $i\in S$ for
which $x_i=−1$.
\end{definition}

\begin{theorem}[Positive and negative coordinate counts]
\lean{FABL.positiveCoordinateCountOn_add_negativeCoordinateCountOn}
\label{FABL.positiveCoordinateCountOn_add_negativeCoordinateCountOn}
\leanok
\uses{FABL.negativeCoordinateCountOn, FABL.positiveCoordinateCountOn}
\difficulty{4}
For any finite set $S$ of coordinates and any sign vector $x$, the number of coordinates
in $S$ where $x$ is $+1$ plus the number where it is $-1$ equals $|S|$.
\end{theorem}

\begin{theorem}[Parity from negative coordinates]
\lean{FABL.parityFunction_eq_neg_one_pow_negativeCoordinateCountOn}
\proofsource{Analysis-of-Boolean-Functions-by-Ryan-ODonnell}{
  pdf-ffbe9c7b3306-p014-b002,
  pdf-ffbe9c7b3306-p023-b001,
  pdf-ffbe9c7b3306-p148-b004
}
\label{FABL.parityFunction_eq_neg_one_pow_negativeCoordinateCountOn}
\leanok
\uses{FABL.Sign, FABL.negativeCoordinateCountOn, FABL.parityFunction}
\difficulty{5}
Let $n$ be a natural number, let $S$ be a subset of the coordinates
$\{0,\ldots,n-1\}$, and let $x\in\{ -1,1\}^n$. The parity function associated to
$S$, evaluated at $x$, equals $(-1)^m$, where $m$ is the number of coordinates in $S$ at
which $x$ equals $-1$.
\end{theorem}

\begin{definition}[Parity count profile]
\lean{FABL.parityCountProfile}
\label{FABL.parityCountProfile}
\leanok
\uses{FABL.Sign}
For a finite coordinate set $S\subseteq\{1,\ldots,n\}$ and an integer $k\ge0$, define
the sign $(-1)^{\max(|S|-k,0)}$.
\end{definition}

\begin{theorem}[Parity from positive-coordinate count]
\lean{FABL.parityFunction_eq_parityCountProfile}
\label{FABL.parityFunction_eq_parityCountProfile}
\leanok
\uses{FABL.negativeCoordinateCountOn, FABL.parityCountProfile, FABL.parityFunction, FABL.parityFunction_eq_neg_one_pow_negativeCoordinateCountOn, FABL.positiveCoordinateCountOn, FABL.positiveCoordinateCountOn_add_negativeCoordinateCountOn}
\difficulty{3}
Let $n$ be a natural number, let $S$ be a subset of the coordinates
$\{0,\ldots,n-1\}$, and let $x\in\{ -1,1\}^n$. If $k$ is the number of coordinates
in $S$ at which $x$ equals $1$, then the parity function associated to $S$, evaluated at
$x$, is $(-1)^{\max(|S|-k,0)}$.
\end{theorem}

\begin{theorem}[Parity as a sum of linear thresholds]
\lean{FABL.parityFunction_isRealSumOfLinearThresholds}
\label{FABL.parityFunction_isRealSumOfLinearThresholds}
\leanok
\uses{FABL.IsRealSumOfLinearThresholds, FABL.isRealSumOfLinearThresholds_of_count_profile, FABL.parityCountProfile, FABL.parityFunction, FABL.parityFunction_eq_parityCountProfile}
\difficulty{1}
Let $n$ be a natural number and let $S$ be a subset of the $n$ coordinates. The parity
function that multiplies the input signs in $S$ is a real constant plus a sum of $2n$
real-valued linear threshold functions on the Boolean cube.
\end{theorem}

\begin{theorem}[Fourier expansion on its support]
\lean{FABL.fourier_expansion_over_support}
\proofsource{Analysis-of-Boolean-Functions-by-Ryan-ODonnell}{
  pdf-ffbe9c7b3306-p022-b002,
  pdf-ffbe9c7b3306-p023-b006
}
\label{FABL.fourier_expansion_over_support}
\leanok
\uses{FABL.fourierCoeff, FABL.fourierSupport, FABL.fourier_expansion, FABL.mem_fourierSupport, FABL.monomial}
\difficulty{4}
Let $n$ be a natural number, let $p$ be a real-valued function on the $n$-dimensional
sign cube, and let $x$ be a point of that cube. The sum, over all coordinate sets with
nonzero Fourier coefficient for $p$, of the Fourier coefficient times the corresponding
sign monomial evaluated at $x$ equals $p(x)$.
\end{theorem}

\begin{theorem}[Threshold Composition Representation]
\lean{FABL.polynomialThresholdRepresentation_isThresholdOfThresholds}
\label{FABL.polynomialThresholdRepresentation_isThresholdOfThresholds}
\leanok
\uses{FABL.BooleanFunction, FABL.BooleanFunction.toReal, FABL.IsPolynomialThresholdRepresentation, FABL.IsThresholdOfThresholds, FABL.fourierCoeff, FABL.fourierSupport, FABL.fourier_expansion_over_support, FABL.monomial, FABL.parityCountProfile, FABL.parityFunction_eq_parityCountProfile, FABL.parityFunction_toReal, FABL.profileDecompositionGates, FABL.profileDecompositionGates_isLinearThreshold, FABL.signValue, FABL.sum_profileDecompositionGates}
\difficulty{6}
Let $n$ be a natural number, let $f:\{−1,1\}^n\to\{−1,1\}$, and let
$p:\{−1,1\}^n\to\mathbb R$. Suppose $f(x)=1$ when $p(x)\geq0$ and $f(x)=-1$ when
$p(x)<0$. Then $f$ is a threshold of $(2n)s$ linear threshold functions, where $s$ is
the number of nonzero Fourier coefficients of $p$ on the Boolean cube.
\end{theorem}

\begin{theorem}[Threshold Composition Size Bound]
\lean{FABL.polynomialThresholdRepresentation_hasThresholdOfThresholdsSizeAtMost}
\label{FABL.polynomialThresholdRepresentation_hasThresholdOfThresholdsSizeAtMost}
\leanok
\uses{FABL.BooleanFunction, FABL.HasThresholdOfThresholdsSizeAtMost, FABL.IsPolynomialThresholdRepresentation, FABL.fourierSupport, FABL.polynomialSparsity, FABL.polynomialThresholdRepresentation_isThresholdOfThresholds}
\difficulty{3}
Let $n$ be a natural number, let $f:\{−1,1\}^n\to\{−1,1\}$, and let
$p:\{−1,1\}^n\to\mathbb R$. Suppose $f(x)=1$ when $p(x)\geq0$ and $f(x)=-1$ when
$p(x)<0$. Then $f$ can be expressed as a threshold of at most $2n\,s$ linear threshold
functions, where $s$ is the number of nonzero Fourier coefficients of $p$ on the Boolean
cube.
\end{theorem}

\begin{theorem}[Sparse polynomial thresholds as threshold compositions]
\lean{FABL.thresholdOfParities_hasThresholdOfThresholdsSizeAtMost}
\label{FABL.thresholdOfParities_hasThresholdOfThresholdsSizeAtMost}
\leanok
\uses{FABL.BooleanFunction, FABL.HasThresholdOfThresholdsSizeAtMost, FABL.IsThresholdOfParities, FABL.polynomialThresholdRepresentation_hasThresholdOfThresholdsSizeAtMost}
\difficulty{4}
Let $f$ be a Boolean function on the sign cube of dimension $n$. If $f$ is represented
as the sign threshold of a real polynomial whose Fourier support has size at most $s$,
then $f$ is a threshold of at most $2ns$ linear threshold functions.
\end{theorem}

\begin{definition}[Complete quadratic bit]
\lean{FABL.completeQuadraticBit}
\statementsource{Analysis-of-Boolean-Functions-by-Ryan-ODonnell}{
  pdf-ffbe9c7b3306-p034-b009,
  pdf-ffbe9c7b3306-p152-b002
}
\label{FABL.completeQuadraticBit}
\leanok
\uses{FABL.𝔽₂}
Let $n$ be a nonnegative integer and let $x=(x_1,\ldots,x_n)\in\mathbb F_2^n$. The
complete quadratic bit of $x$ is $\sum_{1\le i<j\le n}x_i x_j$ in $\mathbb F_2$.
\end{definition}

\begin{definition}[Complete quadratic polar form]
\lean{FABL.completeQuadraticPolar}
\label{FABL.completeQuadraticPolar}
\leanok
\uses{FABL.𝔽₂}
For a natural number $n$ and vectors $x,z\in\mathbb{F}_2^n$, define their complete
quadratic polar form by summing $x_i z_j+z_i x_j$ over all coordinate pairs $i<j$, with
the sum taken in $\mathbb{F}_2$.
\end{definition}

\begin{theorem}[Addition law for the quadratic bit]
\lean{FABL.completeQuadraticBit_add}
\label{FABL.completeQuadraticBit_add}
\leanok
\uses{FABL.completeQuadraticBit, FABL.completeQuadraticPolar}
\difficulty{3}
Let $n$ be a nonnegative integer, and let $x,z\in\mathbb{F}_2^n$, where $\mathbb{F}_2$
is the field with two elements. The complete quadratic bit of $x+z$ equals the sum of
the complete quadratic bits of $x$ and $z$ plus $\sum_{i<j}(x_i z_j+z_i x_j)$, with all
operations in $\mathbb{F}_2$.
\end{theorem}

\begin{definition}[Complete quadratic polar frequency]
\lean{FABL.completeQuadraticPolarFrequency}
\label{FABL.completeQuadraticPolarFrequency}
\leanok
For a natural number $n$ and a vector $z\in\mathbb{F}_2^n$, define a vector whose $i$th
coordinate is the sum in $\mathbb{F}_2$ of all coordinates $z_j$ with $j\ne i$.
\end{definition}

\begin{theorem}[Polar form as a dot product]
\lean{FABL.completeQuadraticPolar_eq_dotProduct}
\label{FABL.completeQuadraticPolar_eq_dotProduct}
\leanok
\uses{FABL.completeQuadraticPolar, FABL.completeQuadraticPolarFrequency, FABL.f₂DotProduct}
\difficulty{5}
Let $n$ be a natural number and let $x,z\in\mathbb F_2^n$. The sum over pairs $i<j$ of
$x_i z_j+z_i x_j$ equals the dot product of $x$ with the vector whose $i$th coordinate
is the sum of all $z_j$ with $j\ne i$.
\end{theorem}

\begin{theorem}[Injectivity of the polar frequency map]
\lean{FABL.completeQuadraticPolarFrequency_eq_zero_of_even}
\label{FABL.completeQuadraticPolarFrequency_eq_zero_of_even}
\leanok
\uses{FABL.completeQuadraticPolarFrequency, FABL.𝔽₂}
\difficulty{5}
Let $n$ be even, and let $z\in\mathbb F_2^n$. If, for each coordinate $i$, the sum in
$\mathbb F_2$ of all coordinates $z_j$ with $j\ne i$ is zero, then $z=0$.
\end{theorem}

\begin{definition}[Real encoding of the complete quadratic]
\lean{FABL.completeQuadratic}
\statementsource{Analysis-of-Boolean-Functions-by-Ryan-ODonnell}{
  pdf-ffbe9c7b3306-p034-b009,
  pdf-ffbe9c7b3306-p152-b002
}
\label{FABL.completeQuadratic}
\leanok
\uses{FABL.completeQuadraticBit, FABL.realSignEncodedFunction}
For each nonnegative integer $n$, define a real-valued function on $\mathbb F_2^n$ by
applying the real sign encoding to the complete quadratic bit.
\end{definition}

\begin{definition}[Boolean complete quadratic function]
\lean{FABL.completeQuadraticBoolean}
\statementsource{Analysis-of-Boolean-Functions-by-Ryan-ODonnell}{
  pdf-ffbe9c7b3306-p034-b009,
  pdf-ffbe9c7b3306-p152-b002
}
\label{FABL.completeQuadraticBoolean}
\leanok
\uses{FABL.BooleanFunction, FABL.binaryCubeSignEquiv, FABL.completeQuadraticBit, FABL.signEncode}
For each nonnegative integer $n$, define a Boolean function on $\{-1,1\}^n$ by
converting each input to its corresponding vector in $\mathbb F_2^n$, evaluating the
complete quadratic bit, and encoding the result as a sign in $\{-1,1\}$.
\end{definition}

\begin{definition}[Symmetric coordinate product]
\lean{FABL.sym2CoordinateProduct}
\label{FABL.sym2CoordinateProduct}
\leanok
\uses{FABL.𝔽₂}
For a natural number $n$ and $x\in\mathbb{F}_2^n$, define a function on unordered pairs
of coordinates by sending the pair $\{i,j\}$ to $x_i x_j\in\mathbb{F}_2$.
\end{definition}

\begin{theorem}[Quadratic bit as a sum over unordered pairs]
\lean{FABL.completeQuadraticBit_eq_sym2_sum}
\label{FABL.completeQuadraticBit_eq_sym2_sum}
\leanok
\uses{FABL.completeQuadraticBit, FABL.sym2CoordinateProduct}
\difficulty{4}
Let $n$ be a nonnegative integer and let $x\in\mathbb{F}_2^n$. The complete quadratic
bit of $x$ equals the sum of $x_i x_j$ over all unordered pairs of distinct coordinates
$\{i,j\}$.
\end{theorem}

\begin{theorem}[Permutation invariance of the quadratic bit]
\lean{FABL.completeQuadraticBit_comp_perm}
\label{FABL.completeQuadraticBit_comp_perm}
\leanok
\uses{FABL.completeQuadraticBit, FABL.completeQuadraticBit_eq_sym2_sum, FABL.sym2CoordinateProduct}
\difficulty{5}
Let $n$ be a nonnegative integer, let $x\in\mathbb{F}_2^n$, and let $\pi$ be a
permutation of the $n$ coordinates. The complete quadratic bit of the vector whose
coordinate $i$ is $x_{\pi(i)}$ equals the complete quadratic bit of $x$.
\end{theorem}

\begin{theorem}[Complete quadratic sign formula]
\lean{FABL.completeQuadratic_apply}
\proofsource{Analysis-of-Boolean-Functions-by-Ryan-ODonnell}{pdf-ffbe9c7b3306-p034-b009}
\label{FABL.completeQuadratic_apply}
\leanok
\uses{FABL.binarySign, FABL.completeQuadratic, FABL.completeQuadraticBit, FABL.realSignEncodedFunction, FABL.signEncodedFunction, FABL.signValue_signEncode_eq_binarySign}
\difficulty{1}
For every natural number $n$ and every $x	ext{in }\F_2^n$, the real sign encoding of
$(-1)^{\text{sum}_{i<j}x_i x_j}$ equals the binary sign character evaluated at
$\text{sum}_{i<j}x_i x_j$. Here the sums and products are taken in $\F_2$.
\end{theorem}

\begin{theorem}[Quadratic Boolean function under sign encoding]
\lean{FABL.completeQuadraticBoolean_binaryCubeSignEquiv}
\proofsource{Analysis-of-Boolean-Functions-by-Ryan-ODonnell}{pdf-ffbe9c7b3306-p034-b009}
\label{FABL.completeQuadraticBoolean_binaryCubeSignEquiv}
\leanok
\uses{FABL.binaryCubeSignEquiv, FABL.completeQuadraticBit, FABL.completeQuadraticBoolean, FABL.signEncode}
\difficulty{1}
For every nonnegative integer $n$ and every $x\in\mathbb{F}_2^n$, evaluate the complete
quadratic Boolean function on the sign vector obtained by encoding each coordinate of
$x$ as $(-1)^{x_i}$. The result is the sign encoding of the complete quadratic bit of
$x$.
\end{theorem}

\begin{theorem}[Symmetry of the quadratic Boolean function]
\lean{FABL.completeQuadraticBoolean_isSymmetric}
\label{FABL.completeQuadraticBoolean_isSymmetric}
\leanok
\uses{FABL.IsSymmetric, FABL.binaryCubeSignEquiv, FABL.binaryCubeSignEquiv_apply, FABL.binarySignEquiv, FABL.completeQuadraticBit_comp_perm, FABL.completeQuadraticBoolean, FABL.permuteInput}
\difficulty{4}
For every nonnegative integer $n$, the complete quadratic Boolean function on $n$ sign
coordinates is invariant under every permutation of its input coordinates.
\end{theorem}

\begin{theorem}[Threshold-of-thresholds size bound]
\lean{FABL.completeQuadraticBoolean_hasThresholdOfThresholdsSizeAtMost_two_mul}
\label{FABL.completeQuadraticBoolean_hasThresholdOfThresholdsSizeAtMost_two_mul}
\leanok
\uses{FABL.HasThresholdOfThresholdsSizeAtMost, FABL.completeQuadraticBoolean, FABL.completeQuadraticBoolean_isSymmetric, FABL.isThresholdOfThresholds_of_isRealSumOfLinearThresholds, FABL.symmetric_hasRealSumOfLinearThresholdsSizeAtMost_two_mul}
\difficulty{3}
For every nonnegative integer $n$, the complete quadratic Boolean function on $n$ sign
coordinates can be expressed as a threshold of at most $2n$ linear threshold functions.
Here each inner function is a linear threshold of the input signs, and the outer
function is a linear threshold of the inner function values.
\end{theorem}

\begin{theorem}[Quadratic bit at zero]
\lean{FABL.completeQuadraticBit_zero}
\label{FABL.completeQuadraticBit_zero}
\leanok
\uses{FABL.completeQuadraticBit}
\difficulty{1}
For every nonnegative integer $n$, the complete quadratic bit of the zero vector in
$\mathbb{F}_2^n$ is $0$.
\end{theorem}

\begin{theorem}[Polar frequency of the zero vector]
\lean{FABL.completeQuadraticPolarFrequency_zero}
\label{FABL.completeQuadraticPolarFrequency_zero}
\leanok
\uses{FABL.completeQuadraticPolarFrequency}
\difficulty{3}
For every natural number $n$, the polar frequency of the zero vector in $\mathbb F_2^n$
is the zero vector.
\end{theorem}

\begin{theorem}[Quadratic translation identity]
\lean{FABL.completeQuadratic_mul_translate}
\label{FABL.completeQuadratic_mul_translate}
\leanok
\uses{FABL.binarySign, FABL.completeQuadratic, FABL.completeQuadraticBit, FABL.completeQuadraticBit_add, FABL.completeQuadraticPolarFrequency, FABL.completeQuadraticPolar_eq_dotProduct, FABL.completeQuadratic_apply, FABL.f₂DotProduct, FABL.vectorWalshCharacter, FABL.vectorWalshCharacter_apply}
\difficulty{4}
For every natural number $n$ and all $y,z	ext{in }\F_2^n$, the product of the complete
quadratic signs at $y$ and $z+y$ equals the binary sign character of the complete
quadratic bit at $z$, multiplied by the Walsh character at $y$ whose frequency has
coordinate $i$ equal to $\text{sum}_{j
e i}z_j$. All vector sums and coordinate sums are over $\F_2$.
\end{theorem}

\begin{theorem}[Self-convolution of a quadratic sign function]
\lean{FABL.convolution_completeQuadratic_self}
\label{FABL.convolution_completeQuadratic_self}
\leanok
\uses{FABL.completeQuadratic, FABL.completeQuadraticPolarFrequency, FABL.completeQuadraticPolarFrequency_eq_zero_of_even, FABL.completeQuadratic_mul_translate, FABL.convolution, FABL.convolution_apply_add, FABL.expect_vectorWalshCharacter}
\difficulty{4}
Let $n$ be an even natural number, and let $z\in (\mathbb{F}_2)^n$. Define
$q(x)=\sum_{0\le i<j<n}x_i x_j$ in $\mathbb{F}_2$, and set $f(x)=(-1)^{q(x)}$. With
convolution normalized by the uniform average on $(\mathbb{F}_2)^n$, the
self-convolution satisfies $(f*f)(z)=1$ if $z=0$, and $(f*f)(z)=0$ otherwise.
\end{theorem}

\begin{theorem}[Fourier coefficient of a singleton indicator]
\lean{FABL.binaryFourierCoeff_singletonIndicator}
\label{FABL.binaryFourierCoeff_singletonIndicator}
\leanok
\uses{FABL.binaryFourierCoeff, FABL.coordinateSum, FABL.χ, FABL.𝔽₂}
\difficulty{4}
Let $n$ be a natural number and let $S$ be any subset of the coordinates of the
$n$-dimensional binary cube. The Fourier coefficient at $S$ of the indicator function of
the zero vector is $2^{-n}$.
\end{theorem}

\begin{theorem}[Squared Fourier coefficient of the complete quadratic]
\lean{FABL.vectorFourierCoeff_completeQuadratic_sq}
\label{FABL.vectorFourierCoeff_completeQuadratic_sq}
\leanok
\uses{FABL.binaryFourierCoeff_convolution, FABL.binaryFourierCoeff_singletonIndicator, FABL.completeQuadratic, FABL.convolution, FABL.convolution_completeQuadratic_self, FABL.f₂Support, FABL.vectorFourierCoeff}
\difficulty{3}
Let $n$ be an even natural number, and let $\gamma\in\mathbb F_2^n$. The square of the
vector Fourier coefficient of the sign encoding of the complete quadratic bit at
$\gamma$ is $2^{-n}$.
\end{theorem}

\begin{theorem}[Fourier coefficient of the complete quadratic function]
\lean{FABL.abs_vectorFourierCoeff_completeQuadratic}
\label{FABL.abs_vectorFourierCoeff_completeQuadratic}
\leanok
\uses{FABL.completeQuadratic, FABL.vectorFourierCoeff, FABL.vectorFourierCoeff_completeQuadratic_sq}
\difficulty{5}
Let $n$ be even, and let $\gamma\in\mathbb F_2^n$. For the real sign-valued function
obtained by encoding the complete quadratic form $\sum_{i<j}x_i x_j$ on $\mathbb F_2^n$,
the absolute value of its Fourier coefficient indexed by the support of $\gamma$ is
$2^{-n/2}$.
\end{theorem}

\begin{theorem}[Real form of the complete quadratic function]
\lean{FABL.completeQuadraticBoolean_toReal}
\label{FABL.completeQuadraticBoolean_toReal}
\leanok
\uses{FABL.BooleanFunction.toReal, FABL.binaryFunctionOnSignCube, FABL.completeQuadratic, FABL.completeQuadraticBoolean, FABL.realSignEncodedFunction, FABL.signEncodedFunction}
\difficulty{1}
For each $n\in\mathbb{N}$, identify the Boolean cube $\mathbb{F}_2^n$ with the sign cube
$\{\pm1\}^n$ coordinatewise, and encode a bit $b$ by the sign $(-1)^b$. The real-valued
function obtained by viewing the sign-encoded complete quadratic Boolean function on
$\{\pm1\}^n$ equals the complete quadratic function transported to that sign cube.
\end{theorem}

\begin{theorem}[Fourier coefficient of a quadratic Boolean function]
\lean{FABL.abs_fourierCoeff_completeQuadraticBoolean}
\label{FABL.abs_fourierCoeff_completeQuadraticBoolean}
\leanok
\uses{FABL.BooleanFunction.toReal, FABL.abs_vectorFourierCoeff_completeQuadratic, FABL.completeQuadraticBoolean, FABL.completeQuadraticBoolean_toReal, FABL.fourierCoeff, FABL.f₂CubeEquivFinset, FABL.f₂Support, FABL.vectorFourierCoeff_eq_fourierCoeff_binaryFunctionOnSignCube}
\difficulty{3}
Let $n$ be an even nonnegative integer and let $S$ be any subset of the $n$ coordinates.
The absolute value of the Fourier coefficient indexed by $S$ of the real-valued complete
quadratic Boolean function on the $n$-dimensional sign cube is $2^{-n/2}$.
\end{theorem}

\begin{theorem}[Sparsity lower bound for a quadratic threshold representation]
\lean{FABL.pow_two_half_le_polynomialSparsity_completeQuadraticBoolean}
\label{FABL.pow_two_half_le_polynomialSparsity_completeQuadraticBoolean}
\leanok
\uses{FABL.IsPolynomialThresholdRepresentation, FABL.abs_fourierCoeff_completeQuadraticBoolean, FABL.completeQuadraticBoolean, FABL.fourierSupport, FABL.one_le_sum_abs_fourierCoeff_of_polynomialThresholdRepresentation, FABL.polynomialSparsity}
\difficulty{5}
Let $n$ be even, and let $p:\{-1,1\}^n\to\mathbb R$ be nonzero. Suppose that, after
identifying each sign vector with its binary vector, the sign of $p(x)$ (with
nonnegative values assigned sign $1$) equals $(-1)^{\sum_{1\le i<j\le n}x_i x_j}$, where
the exponent is computed modulo $2$. Then the number of nonzero Fourier coefficients of
$p$ is at least $2^{n/2}$.
\end{theorem}

\begin{theorem}[Nonconstancy for positive even dimension]
\lean{FABL.completeQuadraticBoolean_ne_one_of_pos_even}
\label{FABL.completeQuadraticBoolean_ne_one_of_pos_even}
\leanok
\uses{FABL.BooleanFunction.toReal, FABL.abs_fourierCoeff_completeQuadraticBoolean, FABL.completeQuadraticBoolean, FABL.expect_monomial, FABL.fourierCoeff}
\difficulty{4}
For every positive even integer $n$, the complete quadratic Boolean function on $n$ sign
coordinates is not the constant function with value $1$.
\end{theorem}

\begin{theorem}[Sparsity of quadratic Boolean thresholds]
\lean{FABL.pow_two_half_le_polynomialSparsity_completeQuadraticBoolean_of_pos}
\label{FABL.pow_two_half_le_polynomialSparsity_completeQuadraticBoolean_of_pos}
\leanok
\uses{FABL.IsPolynomialThresholdRepresentation, FABL.completeQuadraticBoolean, FABL.completeQuadraticBoolean_ne_one_of_pos_even, FABL.polynomialSparsity, FABL.pow_two_half_le_polynomialSparsity_completeQuadraticBoolean, FABL.thresholdSign}
\difficulty{4}
Let $n$ be a positive even integer, and let $p$ be a real-valued function on
$\{-1,1\}^n$ whose threshold sign agrees everywhere with the complete quadratic
Boolean function $x\mapsto (-1)^{\sum_{i<j}x_i x_j}$, using the identification of signs
with bits in $\mathbb F_2$. Then the number of nonzero Fourier coefficients of $p$ is at
least $2^{n/2}$.
\end{theorem}

\begin{theorem}[Threshold parity size lower bound]
\lean{FABL.pow_two_half_le_thresholdOfParitiesSize_completeQuadraticBoolean}
\label{FABL.pow_two_half_le_thresholdOfParitiesSize_completeQuadraticBoolean}
\leanok
\uses{FABL.IsThresholdOfParities, FABL.completeQuadraticBoolean, FABL.pow_two_half_le_polynomialSparsity_completeQuadraticBoolean_of_pos}
\difficulty{4}
Let $n$ be a positive even integer, and let $s$ be a nonnegative integer. If the
complete quadratic Boolean function on $n$ bits has a polynomial threshold
representation with at most $s$ nonzero Fourier coefficients, then $2^{n/2}\le s$.
\end{theorem}
